%
% cover_letter_template_en.tex — English cover letter for English-language
% postings and international/startup environments (DIN 5008-inspired layout).
%
% Same parser-safety recipe as the German template: texgyrepagella with
% Ligatures=NoCommon, four-digit dates, bold subject, collision-proof
% at-token filling.
%
\documentclass[a4paper,11pt]{article}
\usepackage[empty]{fullpage}
\usepackage{geometry}
\usepackage{xcolor}
\usepackage{parskip}
\usepackage[hidelinks]{hyperref}

\geometry{left=2.4cm, right=2cm, top=2.2cm, bottom=2cm}
\pagestyle{empty}
\setlength{\parskip}{0.45em}
\setlength{\parindent}{0pt}

% Parser-safety: ragged right (justified text compresses interword spaces
% below the extractor's x-tolerance and word boundaries are lost in parsing).
\raggedright

% Parser-safety: widen the interword glue (Pagella's 0.25em space extracts
% below the extractor x-tolerance and words merge; see the German template).
\setlength{\spaceskip}{0.34em plus 0.14em minus 0.06em}
\setlength{\xspaceskip}{0.34em plus 0.14em minus 0.06em}

\usepackage{fontspec}
\setmainfont{texgyrepagella}[Ligatures=NoCommon, Extension=.otf, UprightFont=*-regular, BoldFont=*-bold, ItalicFont=*-italic, BoldItalicFont=*-bolditalic]

\newcommand{\clsender}[1]{{\small #1\par}}

\begin{document}

% ---------- sender block (top right) ----------
\hfill
\begin{minipage}[t]{7.2cm}
\raggedleft\clsender{\bfseries Meetkumar Riteshkumar Raval\\ 03046 Cottbus, Germany\\ +49 15567018163 | ravalmeett@gmail.com\\ https://www.linkedin.com/in/meetraval/ | https://github.com/ravalmeet | https://raval.vercel.app/}
\end{minipage}

\vspace{0.9em}

% ---------- company block (left) ----------
Werkstudent*in (m/w/d) IT \& RPA/Automation im KI-Umfeld \\ Germany

\vspace{0.4em}

% ---------- city + date (right, DD.MM.YYYY) ----------
\hfill Cottbus, 29.10.2026

\vspace{1.4em}

% ---------- subject line (bold, exact title) ----------
\noindent\textbf{Application for Werkstudent*in (m/w/d) IT \& RPA/Automation im KI-Umfeld}

\vspace{0.8em}

Dear Hiring Team,

Werkstudent*in IT \& RPA/Automation im KI-Umfeld at Werkstudent*in (m/w/d) IT \& RPA/Automation im KI-Umfeld. My MSc in Artificial Intelligence and experience building automated Python pipelines align directly with this role.

At Soham Infotech I built automated Python data pipelines and backend components, using software engineering best practices and automated workflows. During my internship I created Python-based automated pipelines in Django for data extraction and preprocessing, supporting model training. In the Automated Resume Generation project I developed Python workflows that combine large language models, AI agents and code transformation. My toolkit also includes RPA and KI.

The focus on AI-driven automation of{}fers a compelling environment to apply my expertise in Python and intelligent workflow engineering. I look forward to contributing to innovative RPA solutions within your team.

I am enrolled in the MSc Artificial Intelligence at BTU Cottbus-Senftenberg until October 2027 and available immediately: 20 hours per week during lecture periods and up to 40 hours per week in semester breaks or for an internship, with no fixed end date. My student residence permit is valid until December 2028. I would be glad to discuss my application in a personal or video interview.

\vspace{1.2em}

Kind regards

\vspace{1.6em}

Meetkumar Riteshkumar Raval

\end{document}
